% Abstract (150-250 words typical for CS review journals)
%
% Structure suggestion:
%  - Context/motivation (1-2 sentences)
%  - Problem / gap in current literature
%  - Objective of the review
%  - Methodology (systematic review, databases, years, N studies)
%  - Key findings (3-4 themes)
%  - Implications / contribution
Software testing is undergoing a paradigm shift driven by artificial
intelligence, large-scale automation, and continuous delivery practices.
Despite a growing body of work on next-generation software testing
techniques, the literature remains fragmented across disparate venues,
terminologies, and evaluation methodologies. This paper presents a
systematic literature review (SLR) following the Kitchenham and PRISMA
guidelines to synthesize research published between
\textit{[START\_YEAR]} and \textit{[END\_YEAR]} on AI-driven, autonomous,
and next-generation software testing approaches. We screened
\textit{[N\_INITIAL]} candidate studies from \textit{[DATABASES]},
applying explicit inclusion/exclusion criteria to arrive at
\textit{[N\_FINAL]} primary studies. We classify the selected studies
along four dimensions: testing technique, level of automation, target
application domain, and evaluation rigor. Our synthesis identifies
\textit{[KEY\_THEME\_1]}, \textit{[KEY\_THEME\_2]}, and
\textit{[KEY\_THEME\_3]} as dominant research trends, while highlighting
open challenges in explainability, reproducibility, and industrial
adoption. We conclude with a research agenda for the next generation of
software testing research and practice.
